\scanpage{39}
\mspage{19}
et $j:U-Y$ induit par $\tau_u$, je dis que $H_1(i)=H_1(j)$, ou ce qui revient au $\unsure{\overline m}$, puisque $Y\xrightarrow{\alpha}J_Y^1$ induit un iso. $H_1(\alpha):H_1(Y)-H_1(J_Y^1)$, que $H_1(\alpha)H_1(i)=H_1(\alpha)H_1(j)$, ou encore $H_1(\alpha i)=H_1(\alpha j)$.

Or on a $\alpha j=\tau_u(\alpha i)$ donc \deleted{il suffit de}
\[
H_1(\alpha j)=H_1(\tau_u)H_1(\alpha i)
\]
et on sait que $H_1(\tau_u)=\operatorname{id}$, OK. Cela étant, soit $I=\widehat Y\setminus Y$, $J=\widehat Y\setminus U$, de telle façon que $J\supset I$, et soit $f:\widehat Y\simeq\widehat Y$ déduit de $j:U\subset Y$ en passant aux complétés, donc \inserted{$f(U)=\widehat Y-f(J)\subset Y=\widehat Y-I$} \deleted{\unsure{on a} $f(J)$)} implique $f(J)\supset I$ i.e. $J\supset f^{-1}(I)$.

On voit que le passage de $H_1(U)$ à $H_1(Y)$, par bouchage de trous en $J\setminus I$, a pour noyau la \inserted{somme des} $L_j$ avec $j\in J\setminus I$, et pour la raison analogue le noyau de $H_1(U)-H_1(Y)$ induit par $j$ est la somme des $L_j$ avec $j\in J\setminus f(I)$, d'ailleurs les $L_j$ $(j\in J)$ dans $H_1(U)$ sont $\simeq\widehat{\mathbb Z}$ (resp. $\widehat{\mathbb Z}'$ en car. $p>0$...) et le sous-$\widehat{\mathbb Z}$-module engendré est $\widehat{\mathbb Z}^{J}/\widehat{\mathbb Z}$ (diagonale), donc \deleted{si $\operatorname{card}J\geq3$,} les sous-modules engendrés par deux parties \inserted{$I,I'$} de $J$ distinctes de $J$, ne peuvent être égaux que si $I=I'$, ou $I$ et $I'$ de cardinal $\mu-1$ (où $\mu=\operatorname{card}J$)
\inserted{Mais ici quitte à rapetisser $U$, on peut supposer}
\margin{38}
